
\begin{abstract}
Existing benchmarks for conversational AI agents simulate \textit{single-control} environments, where only the AI agent can use tools to interact with the world, while the user remains a passive information provider. This differs from real-world scenarios like technical support, where users need to actively participate in modifying the state of the (shared) world. In order to address this gap, we introduce \bench, with four key contributions: 
    1) A novel \text{Telecom dual-control domain} modeled as a \decpomdp, where both agent and user make use of tools to act in a shared, dynamic environment that tests both agent coordination and communication,
    2) A \text{compositional task generator} that programmatically creates diverse, verifiable tasks from atomic components, ensuring domain coverage and controlled complexity,
    3) A \text{reliable user simulator} tightly coupled with the environment, whose behavior is constrained by tools and observable states, improving simulation fidelity,
    4) \text{fine-grained analysis of agent performance} through multiple ablations including separating errors arising from reasoning vs communication/coordination.
    In particular, our experiments show significant performance drops when agents shift from no-user to dual-control, highlighting the challenges of guiding users. Overall, \bench provides a controlled testbed for agents that must both reason effectively and guide user actions.\footnote{Code and data are available at: \href{https://github.com/sierra-research/tau2-bench}{https://github.com/sierra-research/tau2-bench}.}
\end{abstract}
